%
% exemplo de artigo usando o modelo RITA.cls
%
\documentclass{RITA}
\usepackage[T1]{fontenc}
%\usepackage[latin1]{inputenc}
\usepackage{url}
\usepackage[brazilian,english]{babel}
\usepackage[utf8]{inputenc}
\usepackage{verbatim}
\usepackage[num]{abntex2cite}	% Citações padrão ABNT

\usepackage[a4paper]{geometry}


%
% preâmbulo
%
\title{Revista RITA: Instruções para Preparação de Documentos em \LaTeX}
\englishtitle{RITA Journal: Instructions for Document Preparing in \LaTeX}
\author{
Maria Lúcia B. Lisbôa
\footnote{Instituto de Informática, UFRGS, Caixa Postal~15064\\
  \texttt{\{llisboa,carla@inf.ufrgs.br\}}}
\and
Carla M. Dal~Sasso Freitas
\footnotemark[1]
\and
Outro Autor
\footnote{Endereço de Outro Autor\\
  \texttt{\{outroautor@email.br\}}}
}
\RITAvolume{24}
\RITAnumber{1}
\RITAyear{2017}
\RITApagerange{1-7}


%
% início do texto
%
\begin{document}
\maketitle

\selectlanguage{brazilian}
\begin{abstract}
As instruções aqui contidas objetivam auxiliar os autores na
preparação de documentos para impressão na Revista RITA.  O
resumo deve ser escrito na mesma língua do texto (Português, Inglês ou
Espanhol) e descreve o conteúdo do texto em cerca de 150-200
palavras. Font Times, tamanho 10, margens laterais reduzidas em 1 cm
de cada lado, com duas linhas em branco antes e depois do
resumo. Formato abstract. Esta é a segunda versão das instruções e
dos formatos e, portanto, sujeita a incorreções e omissões. 
\end{abstract}

\begin{abstractinenglish}
Resumo traduzido para a língua inglesa.
\end{abstractinenglish}

\begin{keywords}
Palavras, chave, em, português.
\end{keywords}

\begin{keywordsenglish}
Keywords, in, english.
\end{keywordsenglish}

\begin{center}                        
    \textit{Data de submissão: 11/06/2016}, \textit{Data de aceite: 02/10/2016}
\end{center}



\section{Seção}
O título de seção deve ser em negrito, numerado automaticamente (1, 1.1,..), Times 12, com 14 espaços antes e 12 após. Este espaçamento pode ser ligeiramente alterado quando houver necessidade de completar uma página.

A preparação do documento deve ser cuidadosa, principalmente quanto a seu conteúdo técnico e correção gramatical. Alertamos os autores que não dispomos de pessoal especializado para fazer correção de textos nas três línguas aceitas (Português, Inglês, Espanhol).

A primeira linha de cada parágrafo possui indentação de 1 cm e as linhas observam espaçamento simples. Cada parágrafo  deve ser separado do seguinte por  6  espaços. Este espaçamento pode ser ligeiramente alterado quando houver necessidade de completar uma página. Estilo text1.

\subsection{Subseção (estilo heading2)}
Cada subseção deve ser em negrito, Times 10, com 12 espaços antes e 8 após. Este espaçamento pode ser ligeiramente alterado quando houver necessidade de completar uma página. Segue-se um parágrafo no estilo text1.

\section{Preparação dos originais}
O texto deve ser preparado usando LaTeX (ou TeX). Por favor veja as instruções contidas no documento RitaLatex.doc.

Ao ler estas instruções, observe na barra de ferramentas o formato correspondente. Use estes mesmos formatos ao preparar seu manuscrito. Por favor, procure seguir as instruções para que a Revista RITA possua uma aparência uniforme.

\subsection{Área de impressão}
A área de impressão é de 13cmX17cm sendo o tamanho da página definido como 16cmX23cm. O texto deve ocupar a linha inteira e as figuras não podem ultrapassar as margens definidas. Procure preencher toda a área de impressão da página de modo não deixar mais do que 2 cm em branco no final de cada página. Faça seu maior esforço, movimentando figuras e texto, para completar cada página.

\subsection{Títulos}
Os títulos, incluindo o título principal, devem iniciar com letra maiúscula, sendo as demais letras minúsculas e todas as demais palavras do título devem iniciar por letra maiúscula. Por exemplo: " Revista RITA: Instruções para Preparação de Documentos".  

\subsection{Cabeçalhos, rodapé e numeração de páginas}
O cabeçalho contém o título do artigo e o rodapé contém a identificação da revista, seu volume e número e a numeração de páginas. Tanto cabeçalho quanto rodapé são espelhados: páginas de numeração ímpar possuem cabeçalho e rodapé alinhados à direita e as de numeração par, à esquerda. Formatos evenhead (caceçalho par), oddhead (cabeçalho ímpar), evenfoot(rodapé par) e oddfoot (rodapé impar). A primeira página não contém cabeçalho; apenas rodapé direito, por iniciar sempre em numeração ímpar.

\subsection{Legendas de tabelas e de figuras}
As tabelas devem ser numeradas a partir de 1 e indicadas por Tabela <num>. <legenda> no topo da mesma, conforme exemplo mostrado na Tabela~\ref{tab:exemplo}. No texto, sempre faça referência a Tabela <num> ( não usar as expressões acima, abaixo, a seguir, etc).  Proceda do mesmo modo com as figuras, lembrando que a legenda da figura fica abaixo da mesma.

\begin{table}
\begin{center}
\caption{Formatos de títulos}
\label{tab:exemplo}
\begin{tabular}{lll}
        \hline
	Item (alinhamento) & Exemplo & Font e estilo\\
	\hline
	Título (esquerdo) & \Large\textbf{Título do Artigo} & 14~pontos, negrito\\
	Seção (esquerdo) & \large\textbf{1~Introdução} & 12~pontos,
	negrito\\
	$\cdots$ & $\cdots$ & $\cdots$\\
	\hline
\end{tabular}
\end{center}
\end{table}

\subsection{Figuras e fotografias}
A impressão da Revista RITA é feita em tons de branco e preto. Portanto, evite fazer uso de fotografias coloridas, a menos que, quando transformadas em tons de cinza seus detalhes continuem visíveis.

As figuras devem ser integradas no texto, centralizadas de acordo com as margens. Para testar a visibilidade dos detalhes de suas figuras, por favor, faça a geração de um arquivo imagem de impressão (postscript) e observem se todos os detalhes estão perfeitamente visíveis e os textos legíveis.  As figuras devem ser numeradas e todas devem ter uma legenda explicativa. Centralize a legenda entre as margens, use font Times 9, e mantenha uma distância de cerca de 6 pontos antes e 12 pontos após (formato legend).

Tenha especial cuidado com figuras feitas diretamente com as ferramentas Word. Permita que elas flutuem sobre o texto, sem âncoras. Se houver necessidade de inserir uma quebra de página antes da figura, certifique-se que a página anterior esteja completa. Se a figura ocupar uma página completa, certifique-se que esta não ultrapasse as margens. Evite colocar figuras e tabelas no formato paisagem.

\subsection{Fórmulas e equações}
Equações e fórmulas deve ser colocadas em uma nova linha,
centralizadas e numeradas consecutivamente para fins de referência,
como pode ser observado na Equação~\ref{eq:exemplo}.

\begin{equation}
x + y = z
\label{eq:exemplo}
\end{equation}

\subsection{Program Code}
Listagens de código de programas devem usar um outro estilo de font, por exemplo, Courier 10, para que recebam destaque.  As listagens não são consideradas figuras, de modo que não necessitam ter legenda. Para fins de referência, as linhas do código podem  ser numeradas. Para numerar a tabela em Word, selecione a primeira célula da tabela, escolha no menu Formatar a opção Numeração (marcadores e numeração) e, na caixa de diálogos exibida selecione a opção Numerada. Ao fazer isto, a partir da primeira linha numerada, a numeração será feita automaticamente a cada troca de linha da tabela (Dica: para inserir uma nova linha, posicione o cursor no final da tabela e pressione Enter).

Por exemplo, o código a seguir mostra uma classe Java, onde a linha 6 inicia um comando que se estende por diversas linhas. Estilo progcode (tabela de duas colunas, com números na coluna 1).

\begin{verbatim}
 1. import java.util.Random;
 2. class Aleatorios {
 3.   public static void main (String[] args) {
 4.     Random qq=new Random();
 5.     for (int k=1;k<10;k++)
 6.       System.out.println("nextInT() da classe Random: " +
 7.         qq.nextInt(100) + "\nMath.random(): " +
 8.         Math.random());
 9.   }
10. }
\end{verbatim}

\subsection{Notas de rodapé}
As notas de rodapé são usadas na primeira página para a identificação dos autores e, ao longo do texto, podem ser usadas para esclarecimentos. 

\subsection{Citações e referências e agradecimentos}
Ao longo do texto, as citações são feitas através de números consecutivos entre colchetes de acordo com a ordem de citação, como em [1],[2],[3]. Ao final, a lista de referências deve ter o nome de Referências (negrito, font Times 12, alinhado à esquerda), sem numeração de seção. Não colocar quebra de página antes. As referências devem ser numeradas, seguindo o texto. Ver exemplo a seguir. 

\subsection{Conferência final}
Ao submeter seu trabalho à Revista RITA, verifique se você possui o seguinte:
Latex: seu arquivo fonte (arquivos TEX para o texto e PS ou EPS para as figuras); arquivos ou macros adicionais que você tenha usado; e o arquivo final DVI. Se possível, prepara um arquivo PDF para a submissão inicial.

Para a submissão inicial enviar apenas o arquivo DVI, ou PS ou PDF.

\section{Contribuição dos Autores}
Ao final do texto, antes da bibliografia, deve estar localizada a sessão Contribuição dos Autores, onde se diz como cada um dos autores contribuiu para o artigo submetido.

\section{Referências}
Esta seção serve somente para esclarecimento da formatação das referências. A seção de referências encontra-se a seguir (sem numeração). Este modelo utiliza o pacote \textit{abntex2cite} para criação e gerenciamento das referências bibliográficas. Para adicionar, remover ou alterar uma referência, basta editar o arquivo "RITAReferencias.bib". As referências que constam no arquivo somente aparecerão no documento se forem citadas ao longo do texto através do comando \textit{cite}. Abaixo é apresentada uma das entradas do arquivo "RITAReferencias.bib", a qual corresponde à referência \cite{Feynman}:

\begin{verbatim}
@Article{Feynman,
         author="Richard P. Feynman", 
         title="The theory of positrons",
         journal="Physical Review",
         address="New York",
         year=1949,
         month=may,
         volume=76,
         number=1,
         pages="749-759"}
\end{verbatim}

O \textit{abntex2cite} aceita diversos tipos de referências. Para inserir uma referência, deve-se escrever o comando correspondente ao tipo desta referência seguido da lista de campos do comando entre chaves, vide exemplo acima. Os comandos mais comuns são "@Article" para artigos, "@Book" para livros, "@Mastersthesis" para teses de mestrado, "@Misc" para websites ou outra categoria que não se encaixa nas demais e assim por diante. Uma lista completa dos tipos de referência aceitos, seus respectivos campos e maiores informações sobre o \textit{abntex2cite} podem ser encontradas em \cite{abntexcite}. A referência utilizada para verificar a corretude do pacote foi \cite{normasABNT}. Após inserida a entrada no arquivo "RITAReferencias.bib" e citada a referência no documento, esta aparecerá automaticamente na seção de referências (vide exemplo abaixo).

\setcounter{secnumdepth}{0}
\bibliography{RITAReferencias}

\end{document}